% DELTA23: author proof and root's independent exact marked-path audit complete.
\subsection{An orbital image of the pentagon with its boundary corrections}

\paragraph{D23.C16 (the five explicitly marked orbital cofaces).}
Use Enriquez's notation
\[
 W_{n,3}=\{(z_1,\ldots,z_n)\in(\mathbb C^\times)^n:
                     z_i^3\ne z_j^3\ (i\ne j)\},
 \qquad \zeta=e^{2\pi i/3}.
\]
There is an algebraic map
\[
 \mathcal O_n:W_{n,3}\longrightarrow M_{0,n+2},\qquad
 (z_1,\ldots,z_n)\longmapsto
       (\zeta z_1,\zeta^2z_1,z_1,z_2,\ldots,z_n).
\tag{D23.C16.1}
\]
All points on the right are distinct. In particular $\mathcal O_3$
induces a homomorphism from the cyclic covering subgroup
$K_{3,3}=\pi_1(W_{3,3})$ to the pure five-point mapping class group.

Fix the following one marking for all five faces. The source base
configurations and connecting paths are in the positive chamber
$0<z_1<z_2<z_3$. Send the fifth target point to infinity by
$M(p)=1/(p-z_3)$, use a projection of angle
$0<\theta<\pi/3$, and straighten the remaining four strands in
reverse label order. Finally conjugate every target word by the
same half twist $a=\sigma_1$. Write
\[
 \begin{gathered}
 a=\sigma_1,\quad b_0=\sigma_2,\quad c=\sigma_3,
 \quad\delta=ab_0,\\
 x=a^2,\quad y=b_0^2,\quad K=c^2,\quad
 Z=(ab_0)^3,\quad V=(b_0c)^3.
 \end{gathered}
\tag{D23.C16.2}
\]
The group is $B_4/\langle z_4\rangle$, where
$z_4=(ab_0c)^4$. Thus $Z$ is the boundary twist $T_{45}$;
$K=T_{34}$ and $V=T_{15}$. In particular
\[
 \delta^3=Z,\qquad [K,x]=[K,V]=1,
       \qquad [Z,x]=[Z,y]=1.
\tag{D23.C16.3}
\]

For the rank-four cyclic source use the actual basis from
Enriquez's Section~8:
\[
 U=x_{01}^3,\qquad Y_j=x_{01}^j x_{12}x_{01}^{-j}\quad(0\le j<3).
\]
The three middle coface maps, followed by $\mathcal O_3$, have
these literal images in the common marking:
\[
\begin{array}{c|cccc}
 &U&Y_0&Y_1&Y_2\\\hline
\widetilde\partial_1&1&Z&1&1\\
\widetilde\partial_2&1&xK^{-1}&V&K(xV)^{-1}\\
\widetilde\partial_3&Z&xZ^{-1}&yZ^{-1}&(xy)^{-1}.
\end{array}
\tag{D23.C16.4}
\]
The exceptional faces are
\[
 \mathcal O_3\partial_0(f)=f(K,Z),\qquad
 \mathcal O_3\partial_4(f)=f(\delta,K).
\tag{D23.C16.5}
\]
The second formula is an extension through the partially labelled
family: rotating the first coordinate through $120$ degrees
permutes the first three target points and gives $\delta$.
Its cube is $Z$, and is not set to one in the five-point group.

Consequently, for every actual full-GT element $h=(1,f)$, put
\[
 b=\Psi_3(h)/2,\qquad C_f=f(x,r)\in F,
 \qquad A^E_f=x^{-3b}C_f^{-1}\in F'.
\]
Here $A^E_f$ is precisely the pure cubic frame of the descent
$\mathcal D_E$ in D23.F22; evaluations of $C_f^{-1}$ or $A^E_f$
at another pair mean evaluation of their marked pure free-group
word. The orbital pentagon identity is
\[
 \boxed{\quad
 f(K,Z)\,A^E_f(x,V)\,K^{-3b}\,f(\delta,K)
                  =A^E_f(x,y)Z^b.
 \quad}
\tag{D23.C16.6}
\]
This is a genuine five-point identity obtained from the original
pentagon, with the two boundary exponents retained.
There is also a specified literal lift to $B_4$, obtained by using
the actual map $(p_1,\ldots,p_5)\mapsto
(1/(p_1-p_5),\ldots,1/(p_4-p_5))$ into the configuration space
of four finite points. In this lift the identity is
\[
 f(K,Z)A^E_f(x,V)K^{-3b}f(\delta,K)
                   =A^E_f(x,y)Z^b z_4^{-b}.
\tag{D23.C16.6L}
\]
In particular a representation with $q(z_4)=-I$ must retain
the scalar $(-I)^b$ on the right. Projective vanishing cannot
be substituted for this literal equation.

\begin{proof}
We first make the base paths explicit. The positive-chamber
connecting paths have no crossings in the chosen target projection.
Indeed put $q=z_1/z_3\in(0,1)$. Then
\[
 \operatorname{Re}(e^{i\theta}M(\zeta^2z_1))
 -\operatorname{Re}(e^{i\theta}M(z_1))
 =\frac{q[3\cos\theta(1+q)-\sqrt3\sin\theta(1-q)]}
        {2z_3(1-q)(1+q+q^2)}>0.
\]
Also $M(z_2)<M(z_1)$ on the real axis, and the imaginary parts
of $M(\zeta z_1),M(\zeta^2z_1)$ give their strict order after the
positive projection. The order is therefore uniformly
$4<3<2<1$ throughout the whole chamber. Thus none of the five
coface identifications acquires a separate hidden conjugator.
The one common conjugation by $a$ changes the observed cluster
rotation $b_0a$ to $ab_0$ and preserves $K$, because $a$ commutes
with $c$.

Use positive lollipop meridians for the base powers
$Z_i=z_i^3$. The generator $x_{01}$ rotates $z_1$ through
$120$ degrees and fixes $z_2$; $x_{12}$ is the lift of a positive
collision loop of $Z_2$ about $Z_1$, along the positive real
connecting arc. Hence $Y_j$ is obtained by rotating the first
coordinate $j$ times, following that collision loop, and rotating
back. This fixes the directions of the sector paths as well as
their endpoints.

The three lifted cofaces of Enriquez's Lemmas~8.2--8.4 are, on
appropriate homotopy-equivalent open sets,
\[
 (z_1,z_2)\longmapsto
 (az_1,z_1,z_2),\quad(z_1,cz_1,z_2),\quad(z_1,z_2,cz_2),
\]
where in this display only $0<a\ll1$ and $c>1$ close to one are
real constants, not braid generators. The collision lollipops
are chosen larger than the indicated splitting distance.

For the first face, the $U$ loop rotates the first four target
points together and is the full four-strand twist, hence trivial
in the sphere quotient. The $Y_0$ loop is the collision of target
points $4,5$, giving $Z$. The other two root-collision divisors
are unmarked regular points in this family, so their meridians
are trivial. This proves the first row. In the ratio
$u=z_2/z_1$, $U=x_{01}^3$ is inertia at $u=\infty$, not at
$u=0$. A meridian at the latter point instead maps to $Z^{-1}$;
confusing those two meridians leaves the value on $f\in F'$
unchanged but gives an incorrect generator table.

For the second face, $U$ again rotates the first four points
and dies. The $Y_0$ loop moves point $5$ around the close pair
$3,4$ without an internal twist. Its value is the twist on
$3,4,5$ divided by $K$, namely $xK^{-1}$. The first rotated
lollipop is the meridian of the remaining orbit point, along
its lower-sector arc. In the common marking its value is
$T_{15}=V$. Both missing peripheral loops $0,\infty$ become
regular, so the three remaining meridians multiply to one;
this gives $Y_2=(xK^{-1}V)^{-1}=K(xV)^{-1}$.

For the third face, $U$ rotates precisely the first three target
points and gives $Z$. The positive-root collision moves the close
pair $4,5$ around point $3$, with no internal twist. Its value is
$xZ^{-1}$. Rotation of the cluster supplies
$Y_j=\delta^j(xZ^{-1})\delta^{-j}$; the braid relations give
$Y_1=yZ^{-1}$ and $Y_2=(xy)^{-1}$. These observations prove
(D23.C16.4). The positive collision arcs in $\partial_0$ give
$K,Z$; the partial cluster rotation and the same collision arc
in $\partial_4$ give $\delta,K$. This proves (D23.C16.5).

There is an independent discrete certificate of the stated braid
identifications. The file
\path{agents/characters/records/orbit_coface_artin_simplifications.json}
records the sixteen crossing words, their one common conjugator,
the simplified words in (D23.C16.4)--(D23.C16.5), and the discarded
power of $z_4$. The faithful Artin action verifies each equality
before quotienting by that specified central power. The numerical
path tracker was used to discover those words; the boundary and
sector-path argument above is the geometric justification.

The path extraction has also been independently certified without
floating-point sampling. In
\path{code/certify_orbital_paths.py}, take source base coordinates
$(z_1,z_2)=(1,4)$, with the extra coordinate $15/2$ when needed.
Use the regular twelve-gon on the unit circle for rotations of
$z_1$: four edges give $x_{01}$, and twelve give $U$.
A collision meridian is a positive diamond of radius $1/20$
times its positive real centre, reached along the positive real
segment and followed by the return segment. For $Y_j$, precede
and follow this lollipop by the prescribed $j$ rotations and
their inverses. These are exactly the positive meridians and
sector arcs specified above. Choose the splitting constants
$a=1/5$ and $c=101/100$, and use the projection factor $1+i/10$.
The diamond about $1$ contains $c$ and $c^{-1}$ in its interior,
but neither $0$ nor the other cube roots of unity; its minimum
distance from $1$ is $1/(20\sqrt2)$. The rotation polygons stay
in the annulus $\cos(\pi/12)\leq|z_1|\leq1$, while their other
source coordinate is $4$. These observations verify that the
source paths and their lifted cofaces stay in the required open
configuration spaces. They are homotopic, with the same sector
markings, to the circular lollipops used to define the cofaces.

All coordinates on a polygon edge are affine functions of its
real parameter with coefficients in $\mathbb Q(\sqrt3,i)$.
After the stated M\"obius normalization, the difference of two
projected real coordinates has a numerator of degree at most
three over $\mathbb Q(\sqrt3)$. Its field norm has degree at
most six over $\mathbb Q$. The certificate isolates every real
root in $(0,1)$ by rational Sturm intervals, separates the two
quadratic-field conjugates by exact signs, and checks the signs
of the derivative and imaginary difference at each retained
root. It also checks the endpoint orders and that no two
crossing intervals overlap. Thus every crossing, its order,
and its sign are proved by rational arithmetic. The resulting
sixteen literal $B_4$ words agree with the recorded words before
any quotient by the centre. The complete root intervals and
polynomials are in \path{records/orbital_paths_exact.json}.

More precisely, in the literal $B_4$ lift the first row of
(D23.C16.4) is $(z_4,Zz_4^{-1},1,1)$ and the second is
\[
       (z_4,xK^{-1}z_4^{-1},Vz_4^{-1},K(xV)^{-1}).
\]
The third row is unchanged. The two generators of the exceptional
zeroth face have images $K,Zz_4^{-1}$, while those of the fourth
face have images $\delta,K$. These are precisely the central
powers retained by the exact Artin certificate.

For the effect on $f$, let its index-three Schreier coefficients
be $(m_0,m_1,m_2)$. D23.F22's integral calculation gives
\[
                    (m_0,m_1,m_2)=(2b,-b,-b),
\]
and the coefficient of $U$ is zero. Filling both fixed punctures
is the map
\[
 \varepsilon_3:(U,Y_0,Y_1,Y_2)\longmapsto
                (1,x,y,(xy)^{-1}),\qquad
 \varepsilon_3(f)=f(r,x)=C_f^{-1}.
\]
Since $K$ commutes with $x,V$, row two of (D23.C16.4) differs
from $C_f^{-1}(x,V)$ by the central factor
$K^{-m_0+m_2}=K^{-3b}$. Since $Z$ commutes with $x,y$,
row three differs from $C_f^{-1}(x,y)$ by
$Z^{-m_0-m_1}=Z^{-b}$. The first row has value $Z^{m_0}=Z^{2b}$.
In particular the half-character correction has been computed
integrally, with no choice of a cube root.

Enriquez's lifted pentagon (51), which is the original full-GT
pentagon inside the covering group, reads
\[
 \partial_0(f)\widetilde\partial_2(f)\partial_4(f)
             =\widetilde\partial_3(f)\widetilde\partial_1(f).
\]
Apply $\mathcal O_3$ and substitute the five formulas. This gives
\[
 f(K,Z)C_f^{-1}(x,V)K^{-3b}f(\delta,K)
                         =C_f^{-1}(x,y)Z^b.
\]
Finally $C_f^{-1}=x^{3b}A^E_f$, and $x$ commutes with both $K$
and $Z$. Cancel the common factor $x^{3b}$ to obtain
(D23.C16.6).
In the literal lift, the first face contributes an additional
$z_4^{-m_0}=z_4^{-2b}$, while the second contributes
$z_4^{-m_0-m_1}=z_4^{-b}$. The central input in the zeroth face
does not change its commutator-word evaluation. Consequently
the same calculation gives (D23.C16.6L). Its total braid exponent
is $-6b$ on each side: the right side has $6b-12b$.
\end{proof}

\paragraph{D23.C16a (the specific cyclotomic action being compared).}
In the orbifold group
\[
 \Gamma=\langle u,v\mid(uv)^2=1\rangle^{\wedge}
           \simeq\widehat{\mathbb Z}\amalg C_2,
\]
Enriquez's Section~13 action, for multiplier one and the canonical
module parameter $g=f$, is
\[
 u\longmapsto u,\qquad
 v\longmapsto f(u,v^2)^{-1}v f(u,v^2).
\]
Filling the two fixed punctures in the order-three cover imposes
$u^3=1$ and gives $B$ with $r=u,s=uv$, so $v^2=x$.
The induced action has $v\mapsto C_fvC_f^{-1}$.
The three-cycle at $(r,srs)$ shows that its action on $s$ is
$f(r,srs)^{-1}s$, exactly the raw action $\psi_f$ of D23.F22.
Thus the descent being compared is a specified cyclotomic orbifold
filling, and not a separately chosen action on an isomorphic group.
The compatibility still sought is with the standard full-GT action
on the five-point periodic frame.

\paragraph{D23.C17 (a precise periodic-factorization residual).}
Let $h=(1,f)$ be an actual full-GT element, and retain the notation
of D23.C16. Let $A_h$ be its original pure cubic frame and put
\[
 D_h=A_h(x,y)^{-1}A^E_f(x,y),\qquad
                g=f(\delta,K).
\]
Then, in the five-point group,
\[
             \boxed{\ (\delta D_hg^{-1}KgD_h^{-1})^3=1.\ }
\tag{D23.C17.1}
\]
Its original signed modular residue is
\[
                  q_3(D_h)=(-I)^{a_3(h)+\Psi_3(h)/2}.
\tag{D23.C17.2}
\]
Therefore excluding the residue $-I$ for this specific correlated
periodic factorization would prove the cubic comparison defect
$e_3=0$ on full GT. This requires less than proving $D_h=1$;
no such exclusion is asserted here.

\begin{proof}
Set $R=f(K,Z)$ and $M=A^E_f(x,V)K^{-3b}$, so $M$ commutes with
$K$. The orbital identity gives
\[
                 (A^E_f)^{-1}R=Z^b g^{-1}M^{-1}.
\]
The standard braid action has
$\Phi(\delta)=A_h\delta A_h^{-1}$ and
$\Phi(K)=RKR^{-1}$. The first identity lifts literally from the
three-strand central quotient: any central discrepancy has braid
exponent zero, while the central generator has exponent six.
Also $(\delta K)^3=z_4=1$ in the sphere quotient. Substitution,
and the fact that $Z$ centralizes $\delta,D_h$, yield
\[
 \Phi(\delta K)
   =A_h Z^b(\delta D_h g^{-1}K gD_h^{-1})Z^{-b}A_h^{-1}.
\]
Cubing proves (D23.C17.1).

The pure frames $A_h,A^E_f$ belong to $F'$. Thus their signed
modulo-three images lie in
$\operatorname{SL}_2(\mathbb F_3)'=Q_8$.
Write the corrected frame as $T_h=(-I)^{a_3(h)}R_0^j$.
Its image in the abelian quotient $C_3$ is that of
$U^{2\rho(h)}$, and $R_0=-UL$ has twice the class of $U$ there;
therefore $j=\rho(h)\bmod3$. D23.F22 gives the same shear for
$\mathcal D_E(h)$ and corrected frame
$(-I)^bR_0^{\rho(h)}$. The shear and order-three factors cancel
in the ratio $q_3(A_h)^{-1}q_3(A^E_f)$, proving (D23.C17.2).
Uniqueness of a periodic frame alone does not make this ratio
trivial: the original and descended actions must still be compared.
\end{proof}
